% Pre-result E1/STRI ReasoningBank Qwen experiment section.
% This file contains design prose only. Populate result language after adjudication.

\subsection{Behavioral propagation through a ReasoningBank representation boundary}
\label{sec:reasoningbank-qwen}

The three studied carriers expose different locations at which representation can
matter. Skill-Pro places an identity-local readiness gate before evolution, so
fragmentation is predicted to prevent otherwise sufficient evidence from reaching
that gate. ACE can merge fragmented evidence, but only a merge that precedes the
consuming operator can repair the boundary. ReasoningBank provides the complementary
positive control: fragments stored within one retrieved case are reunited before
policy consumption, whereas evidence split across cases can remain separated by
top-1 retrieval. We test the latter boundary without changing the semantic items
from which the representations are constructed.

\paragraph{Prospective task-specific memory.}
We construct one frozen memory bank from disjoint SWE-bench Verified source tasks.
The fixed policy runs once on every source task, preserving successful, failed, and
timed-out trajectories. The official ReasoningBank extraction procedure turns each
trajectory into reusable memory items; a deterministic quarter of source tasks is
then reviewed by an independent model for trajectory fidelity, unsupported claims,
test-information leakage, and lesson plausibility. Evaluation tasks retrieve their
top-ranked memory through one fixed embedding and retrieval path. This design makes
memory task-specific through retrieval rather than hand selection, and neither
retrieval relevance nor source success is used to filter the confirmatory sample.

\paragraph{Representations and controls.}
For each evaluation task, treatment A presents the retrieved items as one canonical
case. Treatment D partitions the same underlying items across cases before the
top-1 consumption boundary; it therefore exposes only the first partition to the
policy. Treatment N supplies no ReasoningBank memory and serves as a secondary
uptake diagnostic. Two zero-provider controls localize the intervention: within-case
fragmentation (B) and a case-identifier placebo (E) must serialize to the exact same
complete model request as A, while D must serialize differently. B and E are
structural controls rather than repeated behavioral arms because byte-identical
requests target the same policy distribution.

\paragraph{Population and repeated trials.}
Repositories and tasks are selected by deterministic hash order after historical
exclusions and evaluator qualification, without policy outcomes. Source,
calibration, pilot, and confirmatory tasks are permanently disjoint. The
confirmatory population contains 24 structurally qualified tasks drawn from a
multi-repository design. For each task we execute six planned trials in each of
A, D, and N. These are repeated samples, not retries: every scheduled unit is
attempted once, failed units remain missing, and tasks are never replaced. A frozen
six-round schedule balances arms across early, middle, and late execution positions.

\paragraph{Primary behavioral endpoint.}
The final patch is mapped to an \emph{EditTargetSet}. A Python diff hunk is assigned
to its nearest enclosing function, asynchronous function, or class in the frozen
base-tree AST; unresolved Python hunks use a module/file atom, and non-Python hunks
use a file atom. Jaccard distance between these sets measures edit-target behavior,
not semantic patch equivalence. For task $i$, let $A_i$ and $D_i$ denote its valid
signatures and define
\[
T_i =
2\,\overline d(A_i,D_i)
-\overline d(A_i,A_i')
-\overline d(D_i,D_i').
\]
Thus the cross-arm distance is evaluated relative to within-arm stochastic
dispersion. The global statistic is the mean $T_i$ across tasks with at least four
valid A and four valid D trials. Inference uses 100,000 within-task label
permutations, preserving observed arm sizes, and a 100,000-replicate bootstrap that
resamples tasks rather than trajectories.

\paragraph{Secondary localization.}
The same statistic compares A with N to test whether retrieved memory changes the
behavior distribution at all. Deterministic secondary observables cover the first
action, modified files, exact patch identity, edit targets, tool use, test execution,
trajectory length, and submission state. Official SWE-bench evaluation supplies
per-task repeated-trial resolution proportions; paired task-level analysis keeps
terminal performance distinct from behavioral propagation. Retrieval relevance is
frozen before policy calls, and the preselected high-relevance half is reported only
as secondary sensitivity evidence.

\paragraph{Interpretation.}
A positive A--D result supports propagation from the cross-case boundary to the
distribution of edit-target behavior under this frozen ReasoningBank and Qwen
configuration. A behavioral shift with no detectable terminal difference localizes
attenuation between trajectory and task resolution. If no behavioral separation is
detected, the conclusion remains bounded by the pilot-qualified precision range;
the structural R1 boundary is not reinterpreted as behavioral equivalence.
